% Chapter 4, Topic _Linear Algebra_ Jim Hefferon
%  http://joshua.smcvt.edu/linearalgebra
%  2012-Jun-18
\topic{Metode Chi\`o}
\index{metode Chi\`o|(}

Saat menerapkan Metode Gauss pada matriks yang hanya memuat bilangan bulat,
orang sering ingin agar semua entrinya tetap berupa bilangan bulat.
Untuk menghindari pecahan dalam reduksi matriks ini,
\begin{equation*}
  A=
  \begin{mat}
    2 &1 &1 \\
    3 &4 &-1 \\
    1 &5 &1 
  \end{mat}
\end{equation*}
mereka dapat memulai dengan mengalikan baris-baris bawah dengan~$2$,
\begin{equation*}
  \grstep[2\rho_3]{2\rho_2}
  \begin{mat}
    2 &1 &1 \\
    6 &8 &-2 \\
    2 &10 &2 
  \end{mat}
  \tag{$*$}
\end{equation*}
sehingga eliminasi pada kolom pertama berlangsung sebagai berikut.
\begin{equation*}
  \grstep[-\rho_1+\rho_3]{-3\rho_1+\rho_2}
  \begin{mat}
    2 &1 &1 \\
    0 &5 &-5 \\
    0 &9 &1
  \end{mat}
  \tag{$**$}
\end{equation*}
Pendekatan bilangan bulat sepenuhnya ini lebih mudah untuk perhitungan mental.
Selain itu, penggunaan aritmetika bilangan bulat pada komputer menghindari
sejumlah persoalan rumit yang melibatkan perhitungan titik-mengambang
\cite{Kahan}.
Jadi, pendekatan ini memiliki alasan yang kuat.

Keuntungan lain pendekatan ini adalah kita dapat dengan mudah menerapkan
ekspansi Laplace pada kolom pertama~($**$), lalu memperoleh determinannya
dengan mengingat bahwa hasil tersebut harus dibagi dengan $4$ karena~($*$).

Berikut kasus umum $\nbyn{3}$ bagi pendekatan ini untuk mencari determinan.
Pertama, dengan mengasumsikan $a_{1,1}\neq 0$, kita dapat mengalikan
baris-baris bawah dengan suatu skalar.
\begin{equation*}
  A=
  \begin{mat}
    a_{1,1} &a_{1,2} &a_{1,3}  \\
    a_{2,1} &a_{2,2} &a_{2,3}  \\
    a_{3,1} &a_{3,2} &a_{3,3}  
  \end{mat}
  \grstep[a_{1,1}\rho_3]{a_{1,1}\rho_2}
  \begin{mat}
    a_{1,1}       &a_{1,2}       &a_{1,3}        \\
    a_{2,1}a_{1,1} &a_{2,2}a_{1,1} &a_{2,3}a_{1,1}   \\
    a_{3,1}a_{1,1} &a_{3,2}a_{1,1} &a_{3,3}a_{1,1}   
  \end{mat}
\end{equation*}
Langkah ini mengalikan determinan dengan $a_{1,1}^2$.
Sekarang lakukan eliminasi ke bawah sepanjang kolom pertama.
\begin{equation*}
   \grstep[-a_{3,1}\rho_1+\rho_3]{-a_{2,1}\rho_1+\rho_2}
  \begin{mat}
    a_{1,1}       &a_{1,2}       &a_{1,3}       \\
    0 
       &a_{2,2}a_{1,1}-a_{2,1}a_{1,2} 
       &a_{2,3}a_{1,1}-a_{2,1}a_{1,3}               \\
    0 
       &a_{3,2}a_{1,1}-a_{3,1}a_{1,2} 
       &a_{3,3}a_{1,1}-a_{3,1}a_{1,3}                   
  \end{mat}
\end{equation*}
Misalkan $C$ adalah minor $1,1$.
Menurut ekspansi Laplace, determinan matriks di atas adalah
$a_{1,1}\det(C)$.
Dengan demikian, $a_{1,1}^2\det(A)=a_{1,1}\det(C)$, dan karena
$a_{1,1}\neq 0$, kita memperoleh $\det(A)=\det(C)/a_{1,1}$.

Untuk menangani matriks yang lebih besar, kita harus melihat cara menghitung
entri-entri minornya.
Pola di atas menunjukkan bahwa setiap unsur minor merupakan determinan
$\nbyn{2}$.
Sebagai contoh, entri kiri atas minor,
$a_{2,2}a_{1,1}-a_{2,1}a_{1,2}$, yang merupakan entri $2,2$ pada matriks di
atas, adalah determinan matriks yang dibentuk oleh keempat unsur $A$ berikut.
\begin{equation*}
  \begin{mat}
    \highlight{$a_{1,1}$} &\highlight{$a_{1,2}$} &a_{1,3}  \\
    \highlight{$a_{2,1}$} &\highlight{$a_{2,2}$} &a_{2,3}  \\
    a_{3,1}               &a_{3,2}               &a_{3,3}  
  \end{mat}
\end{equation*}
Demikian pula, entri kiri bawah minor, yaitu entri $3,2$ di atas, merupakan
determinan matriks yang dibentuk oleh keempat unsur berikut.
\begin{equation*}
  \begin{mat}
    \highlight{$a_{1,1}$} &\highlight{$a_{1,2}$} &a_{1,3}  \\
    a_{2,1}               &a_{2,2}               &a_{2,3}  \\
    \highlight{$a_{3,1}$} &\highlight{$a_{3,2}$} &a_{3,3}  
  \end{mat}
\end{equation*}

Jadi, jika $A$ berukuran~$\nbyn{n}$ dengan $n\geq 3$, kita definisikan matriks
Chi\`o $C$ sebagai matriks $\nbyn{(n-1)}$ yang entri $i,j$-nya adalah determinan
\begin{equation*}
  \begin{vmat}
    a_{1,1}  &a_{1,j+1} \\
    a_{i+1,1}    &a_{i+1,j+1}
  \end{vmat}
\end{equation*}
untuk $1\leq i,j\leq n-1$.
\definend{Metode Chi\`o}\index{metode Chi\`o}
untuk mencari determinan $A$ menyatakan bahwa jika $a_{1,1}\neq 0$, maka
$\det(A)=\det(C)/a_{1,1}^{n-2}$.
(Perhatikan bahwa rumus Chi\`o tidak mensyaratkan bilangan bulat; rumus tersebut
juga berlaku bagi bilangan riil.)

Sebagai ilustrasi, kita cari determinan matriks $\nbyn{3}$ berikut.
\begin{equation*} 
  A=
  \begin{mat}
    2 &1 &1 \\
    3 &4 &-1 \\
    1 &5 &1 
  \end{mat}
\end{equation*}
Berikut matriks Chi\`o-nya.
\begin{equation*}
  C=
  \begin{mat}
    \begin{vmat}
      2 &1 \\
      3 &4
    \end{vmat}
   &
   \begin{vmat}
     2 &1 \\
     3 &-1
   \end{vmat}        \\[3ex]      
   \begin{vmat}
     2 &1 \\
     1 &5
   \end{vmat}
   &
   \begin{vmat}
     2 &1 \\
     1 &1
   \end{vmat}
  \end{mat}
  =
  \begin{mat}
    5  &-5  \\
    9  &1
  \end{mat}
\end{equation*}
Rumus bagi matriks $\nbyn{3}$,
$\det(A)=\det(C)/a_{1,1}$, memberikan $\det(A)=(50/2)=25$.

Untuk determinan yang lebih besar, kita harus menjalankan beberapa langkah,
tetapi setiap langkah hanya melibatkan determinan $\nbyn{2}$.
Karena itu, kita sering dapat menghitung determinan dengan menuliskan hanya
sedikit hasil antara.
Sebagai contoh, untuk matriks $\nbyn{4}$ berikut,
\begin{equation*}
  A=
  \begin{mat}
    3  &0  &1  &1  \\
    1  &2  &0  &1  \\
    2  &-1 &0  &3  \\
    1  &0  &0  &1
  \end{mat} 
\end{equation*}
kita dapat mengerjakan setiap perhitungan $\nbyn{2}$ secara mental dan hanya
menuliskan hasil $\nbyn{3}$.
\begin{equation*}
  C_3=
  \begin{mat}
    \begin{vmat}
      3 &0 \\
      1 &2
    \end{vmat}
    &\begin{vmat}
     3 &1 \\
     1 &0 
    \end{vmat}
    &\begin{vmat}
     3 &1 \\
     1 &1
    \end{vmat}                \\[3ex]
    \begin{vmat}
     3 &0 \\
     2 &-1
    \end{vmat}
    &\begin{vmat}
     3 &1 \\
     2 &0
    \end{vmat}
    &\begin{vmat}
     3 &1 \\
     2 &3
    \end{vmat}             \\[3ex]
    \begin{vmat}
     3 &0 \\
     1 &0
    \end{vmat}
    &\begin{vmat}
     3 &1 \\
     1 &0
    \end{vmat}
    &\begin{vmat}
     3 &1 \\
     1 &1
    \end{vmat}
  \end{mat}
  =
  \begin{mat}
    6  &-1  &2 \\
   -3 &-2  &7 \\
    0  &-1  &2
  \end{mat}
\end{equation*}
Perhatikan bahwa determinan matriks ini adalah
$a_{1,1}^{4-2}=3^2$ kali determinan~$A$.

Untuk menyelesaikannya, ulangi proses tersebut.
Berikut matriks Chi\`o dari $C_3$.
\begin{equation*}
  C_2=
  \begin{mat}    
    \begin{vmat}
      6 &-1 \\
     -3 &-2 
    \end{vmat}
    &\begin{vmat}
      6 &2 \\
     -3 &7
    \end{vmat}           \\[3ex]
    \begin{vmat}
      6 &-1 \\
      0 &-1
    \end{vmat}
    &\begin{vmat}
      6 &2 \\
      0 &2
    \end{vmat}          
  \end{mat}
  =
  \begin{mat}
    -15 &48 \\
    -6 &12
  \end{mat}
\end{equation*}
Determinan matriks ini adalah $6$~kali determinan~$C_3$.
Determinan $C_2$ adalah $108$.
Jadi,
$\det(A)=108/(3^2\cdot 6)=2$.

Rumus ekspansi Laplace mengurangi perhitungan determinan~$\nbyn{n}$ menjadi
evaluasi sejumlah determinan~$\nbyn{(n-1)}$.
Rumus Chi\`o juga bersifat rekursif, tetapi mengurangi determinan~$\nbyn{n}$
menjadi satu determinan~$\nbyn{(n-1)}$ yang dihitung dari sejumlah determinan
$\nbyn{2}$.
Namun, bagi matriks besar, Metode Gauss lebih baik daripada keduanya; sebagai
contoh, metode tersebut memerlukan kira-kira setengah dari jumlah operasi yang
diperlukan Metode Chi\`o~\cite{FullerLogan}.


\begin{exercises}
  \item 
    Gunakan Metode Chi\`o untuk mencari setiap determinan.
    \begin{exparts*}
      \partsitem 
        $
        \begin{vmat}
          1  &2  &3  \\
          4  &5  &6  \\
          7  &8  &9  
        \end{vmat}
        $
      \partsitem 
        $
        \begin{vmat}
          2  &1  &4  &0  \\
          0  &1  &4  &0  \\
          1  &1  &1  &1  \\
          0  &2  &1  &1  
        \end{vmat}
        $
    \end{exparts*}
    \begin{answer}
      \begin{exparts*}
        \partsitem Matriks Chi\`o-nya adalah
        \begin{equation*}
          C=
          \begin{mat}
            -3 &-6 \\
            -6 &-12
          \end{mat}
        \end{equation*}
        dan determinannya adalah $0$.
        \partsitem Mulailah dengan
        \begin{equation*}
          C_3=
          \begin{mat}
            2 &8  &0 \\
            1 &-2 &2 \\
            4 &2  &2
          \end{mat}
        \end{equation*}
        lalu langkah berikutnya memberikan
        \begin{equation*}
          C_2=
          \begin{mat}
            -12 &4 \\
            -28 &4
          \end{mat}
        \end{equation*}
        dengan determinan $\det(C_2)=64$.
        Dengan demikian, determinan matriks semula adalah
        $64/(2^2\cdot 2^1)=8$.
      \end{exparts*} 
    \end{answer}
  \item Bagaimana jika $a_{1,1}$ bernilai nol?
    \begin{answer}
      Jika matriks memiliki entri tak nol $a_{i,j}$, tukarkan baris~$i$ dengan
      baris pertama dan kolom~$j$ dengan kolom pertama untuk membentuk matriks
      $A'=(a'_{p,q})$ dengan $a'_{1,1}\neq 0$.
      Setiap pertukaran membalik tanda determinan; setelah tanda ini dicatat,
      terapkan konstruksi biasa pada $A'$.
      Entri $c_{p,q}$ matriks Chi\`o bagi $A'$ adalah
      \begin{equation*}
        \begin{vmat}
          a'_{1,1}    &a'_{1,q+1} \\
          a'_{p+1,1}  &a'_{p+1,q+1}
        \end{vmat}
      \end{equation*}
      untuk $1\leq p,q\leq n-1$.
      Jika semua entri $A$ bernilai nol, maka tentu saja $\det(A)=0$.
    \end{answer}
  \item \definend{Aturan Sarrus}\index{Sarrus, Aturan}\index{Aturan Sarrus}
    merupakan alat bantu ingatan yang banyak dipelajari orang untuk rumus
    determinan $\nbyn{3}$.
    Di sebelah kanan matriks, salin dua kolom pertama.
    \begin{equation*}
      \begin{array}{ccc|cc}
        a &b &c &a &b \\
        d &e &f &d &e \\
        g &h &i &g &h
      \end{array}
    \end{equation*}
    Determinannya kemudian merupakan jumlah tiga diagonal dari kiri atas ke
    kanan bawah dikurangi jumlah tiga diagonal dari kiri bawah ke kanan atas,
    $aei+bfg+cdh-gec-hfa-idb$.
   Hitung banyak operasi yang terlibat dalam rumus Sarrus dan rumus Chi\`o.
   \begin{answer}
     Rumus Sarrus menggunakan $12$~perkalian dan $5$ penjumlahan (dengan
     pengurangan dihitung bersama penjumlahan).
     Rumus Chi\`o menggunakan dua perkalian dan satu penjumlahan
     (yang sebenarnya berupa pengurangan) bagi masing-masing dari empat
     determinan $\nbyn{2}$, lalu dua perkalian dan satu penjumlahan lagi bagi
     determinan Chi\`o $\nbyn{2}$, serta pembagian terakhir dengan $a_{1,1}$.
     Jumlahnya sebelas perkalian atau pembagian dan lima penjumlahan atau
     pengurangan. Jadi, Chi\`o unggul.
   \end{answer}
  \item Buktikan rumus Chi\`o.
    \begin{answer}
      Tinjau matriks $\nbyn{n}$.
      \begin{equation*}
        A=
        \begin{mat}
           a_{1,1}  &a_{1,2}   &\cdots &a_{1,n-1}  &a_{1,n} \\
           a_{2,1}  &a_{2,2}   &\cdots &a_{2,n-1}  &a_{2,n} \\
                  &\vdots                         \\
           a_{n-1,1} &a_{n-1,2} &\cdots &a_{n-1,n-1} &a_{n-1,n}  \\ 
           a_{n,1}  &a_{n,2}   &\cdots &a_{n,n-1}  &a_{n,n} 
        \end{mat}
      \end{equation*}
      Kalikan setiap baris selain baris pertama dengan~$a_{1,1}$.
      \begin{equation*}
        \grstep[a_{1,1}\rho_3 \\ \vdots \\ a_{1,1}\rho_{n}]{a_{1,1}\rho_2}  
        \begin{mat}
          a_{1,1}         &a_{1,2}        &\cdots &a_{1,n-1}        &a_{1,n} \\
          a_{2,1}a_{1,1}   &a_{2,2}a_{1,1}  &\cdots &a_{2,n-1}a_{1,1}  &a_{2,n}a_{1,1} \\
                        &\vdots                         \\
          a_{n-1,1}a_{1,1} &a_{n-1,2}a_{1,1} &\cdots &a_{n-1,n-1}a_{1,1} &a_{n-1,n}a_{1,1}\\
          a_{n,1}a_{1,1}  &a_{n,2}a_{1,1}   &\cdots &a_{n,n-1}a_{1,1}  &a_{n,n}a_{1,1} 
        \end{mat}
      \end{equation*}
      Langkah tersebut mengalikan determinan dengan faktor~$a_{1,1}^{n-1}$.

      Selanjutnya, lakukan operasi baris $-a_{i,1}\rho_1+\rho_i$ pada setiap
      baris~$i>1$.
      Operasi-operasi baris ini tidak mengubah determinan.
      \begin{equation*}
        \grstep[-a_{3,1}\rho_1+\rho_3 \\ \vdotswithin{-a_{3,1}\rho_1+\rho_3} \\ -a_{n,1}\rho_1+\rho_n]{-a_{2,1}\rho_1+\rho_2}         
      \end{equation*}
      Hasilnya adalah matriks~$B$ yang baris pertamanya tidak berubah, kolom
      pertamanya seluruhnya nol (kecuali entri $1,1$, yaitu $a_{1,1}$), dan
      entri-entri sisanya adalah sebagai berikut.
      \begin{equation*}
        \begin{mat}
              a_{1,2}  
              &\cdots 
              &a_{1,n-1}  
              &a_{1,n} 
              \\
              a_{2,2}a_{1,1}-a_{2,1}a_{1,2}  
              &\cdots 
              &a_{2,n-1}a_{1,1}-a_{2,1}a_{1,n-1}
              &a_{2,n}a_{1,1}-a_{2,1}a_{1,n}
              \\
              \vdots                         \\
             a_{n,2}a_{1,1}-a_{n,1}a_{1,2} 
             &\cdots 
             &a_{n,n-1}a_{1,1}-a_{n,1}a_{1,n-1} 
             &a_{n,n}a_{1,1}-a_{n,1}a_{1,n}  
        \end{mat}
      \end{equation*}
       Determinan matriks $\nbyn{n}$ ini, yaitu $B$, adalah
       $a_{1,1}^{n-1}$~kali determinan~$A$.

     Nyatakan minor $1,1$ dari matriks $B$ dengan~$C$, yaitu submatriks yang
     terdiri atas baris~$2$ sampai~$n$ dan kolom~$2$ sampai~$n$.
     Ekspansi Laplace sepanjang kolom pertama $B$ memberikan
     $\det(B) = (-1)^{1+1}a_{1,1}\det(C)$.

      Jika $a_{1,1}\neq 0$, menyamakan kedua ekspresi bagi $\det(B)$ dan
      mencoret faktor yang sama memberikan
      $\det(A)=\det(C)/a_{1,1}^{n-2}$.
    \end{answer}
\end{exercises}

\announcecomputercode
Kode berikut mengimplementasikan Metode Chi\`o.
Kode ini ditulis dalam bahasa komputer Python.
%  (to make it more readable 
% the code avoids some Python facilities).
% Note the recursive call in the final line of 
% \lstinline!chio_det!. 
\begin{lstlisting}
#!/usr/bin/env python3
# chio.py
#  Hitung determinan menggunakan metode Chio.
# Jim Hefferon; Domain Publik
# Hanya untuk demonstrasi; misalnya, tidak menangani kasus M[0][0]=0

def det_two(a,b,c,d):
    """Kembalikan determinan matriks 2x2 [[a,b], [c,d]]"""
    return a*d-b*c

def chio_mat(M):
    """Kembalikan matriks Chio sebagai daftar baris
        M  matriks nxn, berupa daftar baris"""
    dim=len(M)
    C=[]
    for row in range(1,dim):
        C.append([])
        for col in range(1,dim):  
            C[-1].append(det_two(M[0][0], M[0][col], M[row][0], M[row][col]))
    return C

def chio_det(M,show=None):
    """Cari determinan M dengan metode Chio
        M  matriks mxm, berupa daftar baris"""
    dim=len(M)
    key_elet=M[0][0]
    if dim==1:
        return key_elet
    return chio_det(chio_mat(M))//(key_elet**(dim-2))


if __name__=='__main__':
    M=[[2,1,1], [3,4,-1], [1,5,1]]
    print("M=", M)
    print("Determinan =", chio_det(M))
\end{lstlisting}
Berikut hasil pemanggilan program dari baris perintah.
\begin{lstlisting}
$ python3 chio.py
M= [[2, 1, 1], [3, 4, -1], [1, 5, 1]]
Determinan = 25
\end{lstlisting}
\index{metode Chi\`o|)}
\endinput
